% Generated draft: requires mathematical and source review.
\documentclass[11pt]{article}
\usepackage[margin=25mm]{geometry}
\usepackage{fontspec}
\setmainfont{Latin Modern Roman}
\usepackage{amsmath,amssymb,mathtools,unicode-math}
\setmathfont{Latin Modern Math}
\usepackage{enumitem,needspace,xurl,hyperref}
\setlength{\parindent}{0pt}
\setlength{\parskip}{5pt}
\setlength{\emergencystretch}{3em}
\newcommand{\N}{\mathbb{N}}
\newcommand{\Z}{\mathbb{Z}}
\newcommand{\Q}{\mathbb{Q}}
\newcommand{\R}{\mathbb{R}}
\newcommand{\C}{\mathbb{C}}
\newcommand{\F}{\mathbb{F}}
\newcommand{\A}{\mathbb{A}}
\newcommand{\PP}{\mathbb{P}}
\newcommand{\E}{\mathbb{E}}
\newcommand{\eps}{\varepsilon}
\newcommand{\dd}{\,\mathrm{d}}
\begin{document}
\section*{UnsolvedMath ID: 20000383}
\textbf{Source problem code:} catalogue entry
\subsection{Definitions and mathematical statement}
A cyclic triple cover of \(\PP^1\) means a degree-three finite morphism \(\pi:C\to\PP^1\) from a smooth projective curve, whose associated function-field extension is Galois with cyclic group of order three. The question is whether the two properties denoted \((A)\) and \((B)\) in the workshop discussion can be recovered for such covers when the family is ordered by the genus \(g(C)\). The supplied record does not define \((A)\), \((B)\), the ground field, or the equivalence and weighting conventions for the family, so it does not determine two precise limit assertions.

The supplied research report addresses the following narrower setting: fix a finite field \(\F_q\) with \(q\equiv1\pmod3\), and consider the full genus-\(g\) moduli family with reciprocal-automorphism weights. It attributes to Bucur, David, Feigon, Kaplan, Lal\'in, Ozman and Wood a theorem asserting asymptotic independence of prescribed local ramified, split and inert behaviors with explicit masses, Chebotarev proportions conditional on being unramified, and an independent identically distributed limit for the point count as \(g\to\infty\). The local masses, the limiting point-count distribution, and the precise automorphism/moduli convention are not given in this record. These reported results cannot be assigned to the undefined labels \((A)\) and \((B)\).

A separate, explicitly stated candidate proposition concerns only labelled geometric inertia with a chosen generator of the cyclic group. For an integer \(r\ge2\), put
\[
 T_r=\{(a_1,\ldots,a_r)\in\{1,2\}^r:\ a_1+\cdots+a_r\equiv0\pmod3\}.
\]
Give \(T_r\) uniform probability. The reported proposition is
\[
 |T_r|=\frac{2^r+2(-1)^r}{3},\qquad
 \Pr(a_i=1)=\Pr(a_i=2)=\frac12\quad(1\le i\le r),
\]
and, for \(1\le k\le r\),
\[
 d_{\mathrm{TV}}\bigl(\mathcal L(a_1,\ldots,a_k),U_k\bigr)
 \le \min\left\{1,\frac{2^k+1}{2^r-2}\right\}.
\]
Here \(U_k\) is the uniform product law on \(\{1,2\}^k\), and \(d_{\mathrm{TV}}(\mu,\nu)=\tfrac12\sum_x|\mu(x)-\nu(x)|\). In tame cyclic triple Hurwitz data with all \(r\) geometric branch points fully ramified, Riemann--Hurwitz gives \(r=g+2\). Uniform labelled, generator-marked inertia tuples are the specific probability space of this proposition; the assertion does not by itself identify this law with the arithmetic distribution of unlabelled covers. The restrictions \(r\ge2\) and \(1\le k\le r\) make the displayed probability law and denominator meaningful.
\subsection{Short English statement}
The source asks for two unnamed properties of cyclic triple covers when covers are counted by genus. The supplied research report describes finite-field local-statistics results in a tame weighted family and a finite-genus correction for labelled inertia tuples, but the meanings of \((A)\) and \((B)\) remain missing.
\subsection{Sources}
\begin{itemize}
\item Local UnsolvedMath registry, permanent ID \(20000383\), source code \texttt{catalogue entry}; AIM algebraic number theory corpus, \texttt{aim-algebraic-number-theory-notes.json}, notes[34]. Attribution: Ulam AI / UnsolvedMath Contributors, CC BY 4.0.
\item AIM workshop \emph{Arithmetic statistics over finite fields and function fields}, Question 5: \url{https://aimath.org/pastworkshops/arithstatffieldproblems.pdf}. Original remarks and literature fields are empty. The verification field labels the statement exact but ends with ``The exact canonical record is:'' and supplies no further verification detail.
\item The supplied research summary names a theorem of Bucur, David, Feigon, Kaplan, Lal\'in, Ozman and Wood, but gives no title, date or precise bibliographic citation. No such details are invented here.
\item Registry assessment dated August 14, 2026: \texttt{partially\_solved}, \texttt{PARTIAL-PROGRESS}; attempt 1 labelled partial result, rigor proved, novelty confidence low. The report and its candidate proposition are machine-generated research aids, not peer reviewed, requiring independent verification before citation.
\end{itemize}
\end{document}
